\section{Conclusion}

I develop an endogenous-growth model in which AI capability makes compute more
productive both in current production and in AI research. This creates two
feedback loops. The technological loop runs from capability to more effective
research compute and back to capability. The economic loop runs from capability
to output, accumulation, and research expenditure. Following the logic of
\citet{davidsonetal2026}, the strength of these loops depends not only on
technology but also on equilibrium prices and allocations.

The first result concerns the developer's research problem. On every finite
horizon, \(\eta<\alpha\) makes the objective coercive in total research
expenditure and keeps its value finite. When both CES technologies permit gross
substitution, \(\eta>\alpha\) instead allows research-compute expansions that
make the value unbounded. The benchmark \(\eta=0.20<\alpha=0.33\) is therefore a
deliberate regularity restriction, not an empirical estimate. Coercivity is
necessary for a well-behaved problem, but it does not by itself prove existence
of a global optimum or an infinite-horizon equilibrium.

The final-production elasticity then organizes the macroeconomic results. When
labor and AI production services are complements, labor remains a bottleneck and
capability creates no per-capita growth beyond exogenous labor productivity. In
the reported case with zero exogenous productivity growth, output per person
converges to a constant even though capability continues to improve. At unit
elasticity, the automated-research benchmark admits an equilibrium
balanced-growth path under explicit restrictions, and the calibration has the
local saddle-path structure needed for transition equilibrium from nearby
stocks. In the human-research extension, the corresponding growth restrictions
remain necessary rather than sufficient. The reciprocal of the feedback
denominator is the scalar counterpart of the feedback multiplier in
\citet{davidsonetal2026}. When labor and AI production services
are gross substitutes and AI production services dominate the final-production composite,
\(Y/K\) and the return to capital become unbounded, so no finite-rate
balanced-growth path can describe that branch.

These regimes also have sharply different implications for factor payments and
resource use. The outer Cobb--Douglas technology fixes the final-good firm's
gross capital payment
at \(\alpha Y\), including depreciation. What changes with \(\sigma_{XL}\) is
the division of the remaining output between production labor and purchases of
AI production services. Those purchases are not developer profit: they finance inference
compute, human research, and research compute before the residual \(\Pi\) is
distributed. On the conditional gross-substitutes branch, aggregate labor
income vanishes relative to output while compute uses and distributed developer
profit expand. Yet the real wage can still rise because output grows faster
than labor income. This is a decomposition of factor payments and resource
uses, not a prediction about interpersonal inequality: the representative
household owns capital and the developer, and the model does not specify who
ultimately receives payments for upstream compute.

The broad displacement mechanism is known. Task models distinguish automation
from the creation of new human tasks \citep{acemoglurestrepo2019}, and
\citet{restrepo2025agi} shows that labor's income share can converge to zero once
compute performs all bottleneck work. My contribution is to connect that
possibility to two explicit substitution margins in a decentralized Ramsey
economy and to separate physical-capital payments, production and research
wages, inference and research compute, and distributed developer profit along
the transition.

The timing result that survives the equilibrium-admission requirement is
one-sided. An equilibrium with \(\sigma_{XL}=0.99\), started from the same
predetermined stocks as the unit-elastic balanced-growth equilibrium, remains
close to it for a long time before level and distributional differences become
material. Analytically, nearby dated solutions separate on a time scale of
order \(1/|\sigma_{XL}-1|\). I do not report the corresponding
\(\sigma_{XL}>1\) trajectory. Bounded capability cannot coexist with a bounded
interest rate, while the regular AI-dominated continuation ends at a finite
singular date. The proof does not yet exclude bounded capability with
unbounded interest-rate spikes or sufficiently irregular unbounded capability.

The numerical evidence has a precise but limited interpretation. Every
reported trajectory is a numerical approximation to an equilibrium admitted by
the dated equations, a regime-specific long-run continuation, sufficient
optimality conditions, and both transversality conditions. These checks do not constitute global
existence or uniqueness theorems. Gross-substitutes free-boundary calculations
satisfy the dated necessary conditions but fail the equilibrium-admission rule;
they therefore generate no reported trajectory.

Three extensions would be especially useful. An upstream compute sector would
determine how compute expenditure is divided among hardware, datacenters,
energy, labor, and rents. Separate software, hardware, and general-productivity
stocks would bring the richer feedback network of \citet{davidsonetal2026} into
the equilibrium allocation problem. Finally, allowing human researchers to
provide supervision as well as ideas would connect research automation to
competition, control, and AI governance. Each extension changes an economic
mechanism rather than merely adding notation, and therefore belongs in a model
designed expressly for that question.
